\cabstract{
在学习王老师的大数据原理与技术课程\cite{lecture01}后，在此报告中完成相应的课后作业。
本作业围绕混合推荐系统设计作业展开。

\textbf{任务要求：}
\begin{itemize}
    \item 基于 MovieLens 数据集
    \item 实现协同过滤（用户/物品相似度）与基于内容的推荐（电影标签），并比较两者的召回率；
    \item 设计加权混合策略或其他策略以提升推荐多样性。
\end{itemize}

\textbf{提交要求：}
\begin{itemize}
    \item 提交内容：代码 + 简单报告（至少 2 页）
    \item 命名格式：学号+姓名+作业7.zip
    \item 提交至邮箱：BDSysu2026Spring@163.com
    \item 截止时间：布置作业之后的下一个周日 24:00
    \item 严禁抄袭：任何一次抄袭，平时成绩按零分处理
\end{itemize}

本报告基于 MovieLens（ml-latest-small）数据集，设计并实现了一个可解释的混合电影推荐系统。系统包括：
1) 基于用户/物品相似度的协同过滤（User-CF 与 Item-CF）；
2) 基于电影类型（genres）的内容推荐（TF-IDF + 余弦相似度）；
3) 将上述子系统以加权混合策略融合，并在冷启动场景引入流行度兜底与多样性惩罚机制。

报告中详细记录了数据预处理、推荐系统设计、候选召回与融合流程、实验结果与案例分析等内容。实验结果表明，混合策略在多数示例用户上能提升 Recall@20，且在历史记录稀少的用户上通过热度补偿提高覆盖率；本次作业实现了“召回 + 融合 + 重排”流水线的工程，也进行了完备的实验效果分析。
}
\ckeywords{混合推荐系统设计，协同过滤，基于内容的推荐，召回率，多样性提升}
\ekeywords{}


